% !TEX root = main04_en_noimage.tex
\documentclass[reprint,amsmath,amssymb,aps,prd]{revtex4-2}

\usepackage{dcolumn}
\usepackage{bm}
\usepackage{hyperref}

\begin{document}

\preprint{APS/123-QED}

\title{Mechanics of Spatial Vibration II: Interaction of Macroscopic Gravity with Microscopic Vibration and a Unified Dark Fluid Model}

\author{Kwang Yong Yoo}
\email{khanyong.yoo@gmail.com}
\affiliation{Independent Researcher \\ 
\textit{AI-Assisted Research and Simulation Development Support: Gemini} \\
(Interactive simulations and supplementary materials are officially archived under DOI: \href{https://doi.org/10.5281/zenodo.21233252}{10.5281/zenodo.21233252})}

\date{\today}

\begin{abstract}
To address the cosmological challenges associated with anomalous galactic rotation curves and the accelerating expansion of the universe, this study explores a geomechanical alternative to standard dark matter and dark energy paradigms. Building upon the `Geometrical Fluctuation of Space' hypothesis introduced in Part I, this paper proposes an effective unified model by incorporating a phenomenological `Spatial Vibration Tensor' ($\tilde{V}_{\mu\nu}$) into the Einstein Field Equations. 

To preserve the equivalence principle and local energy-momentum conservation, spatial vibration is modeled as high-frequency perturbations to the background metric, rigorously averaged via the Isaacson limit ($\tilde{V}_{\mu\nu} \propto \langle \nabla_\mu h_{\alpha\beta}^{vib} \nabla_\nu h^{\alpha\beta}_{vib} \rangle$). We propose a `Unified Dark Fluid' behavior employing a heuristic chameleon-like transition function: in dense galactic halos, constructive interference traps kinetic modes, emerging as an effective mass ($w \approx 0$) that generates surplus gravity (phenomenological Dark Matter). Conversely, in vast cosmic voids, the unsuppressed potential tension of the vacuum acts as a repulsive vibrational radiation pressure ($w \to -1$), driving accelerated expansion (Dark Energy). Furthermore, micro-perturbations in gravitational lensing are rigorously modeled as stochastic geodesic deviations induced by high-frequency fluctuations ($h^{vib}$), predicting `Achromatic Optical Blurring'. Chronological illusions induced by spatial junctions are proposed as `Geometrical Path-Shortening' to address discrepancies in early galaxy observations.
\end{abstract}

\maketitle

\section{Introduction}
In Paper I of this series, a framework was proposed wherein probabilistic wave phenomena in the microscopic realm are modeled as arising from the geometric fluctuations of the spatial background. It was also established that dense macroscopic masses dynamically suppress this local guidance force due to their inertia, leading to an effective convergence toward Newtonian classical mechanics.

However, when scaling up to cosmological dimensions, the dynamic interplay between mass inertia and spatial volume undergoes a dramatic inversion. The immense spatial vibrational energy accumulated within the vast volumes of deep space seizes the opportunity of diluted mass density to once again overpower the macroscopic inertia of celestial bodies---a phenomenon defined here as \textbf{`Scale Inversion'}.

From this macroscopic perspective, this paper investigates the mechanical interactions between the macroscopic spacetime curvature described by General Relativity and the microscopic geometrical fluctuation of space. By formulating a feedback loop where spacetime curvature distorts spatial vibrations, which in turn dynamically contribute to the effective energy-momentum tensor, this study attempts to reinterpret the phenomena attributed to dark matter and dark energy as consequences of the mechanical phase transitions of empty space itself.

\section{Coupled Equations of Vibration and Gravity}
The standard Einstein Field Equations typically assume the vacuum state to be a smooth continuum \cite{Einstein1916}. However, according to this study's axioms, the vacuum possesses geometrical amplitudes fluctuating at microscopic scales.

\subsection{Spatial Vibration Tensor and Isaacson Averaging}
To describe the kinetic energy and geometrical repulsive radiation pressure inherently emitted by the vibrating empty space, a phenomenological \textbf{`Spatial Vibration Tensor' ($\tilde{V}_{\mu\nu}$)} is introduced as an additional source term:
\begin{equation}
G_{\mu\nu} = \frac{8\pi G}{c^4} \left( T_{\mu\nu} + \tilde{V}_{\mu\nu} \right)
\end{equation}
The cosmological constant $\Lambda$ is entirely discarded. To strictly satisfy the Bianchi Identity ($\nabla^\mu G_{\mu\nu} = 0$), the vibration tensor is mandated to conserve energy-momentum locally. 

Spatial vibration is treated as high-frequency metric perturbations ($g_{\mu\nu} = \bar{g}_{\mu\nu} + h_{\mu\nu}^{vib}$). Since the linear spatial average of a pure wave is exactly zero ($\langle h_{\mu\nu}^{vib} \rangle = 0$), the effective spatial vibration tensor $\tilde{V}_{\mu\nu}$ is mathematically constructed analogously to the Isaacson high-frequency limit \cite{Isaacson1968}, ensuring strict covariance:
\begin{equation}
\tilde{V}_{\mu\nu} \propto \frac{c^4}{G} \langle \nabla_\mu h_{\alpha\beta}^{vib} \nabla_\nu h^{\alpha\beta}_{vib} \rangle
\end{equation}

\subsection{Dynamical Feedback Loop}
The extended equation reveals a nonlinear feedback loop. As a massive object curves the background spacetime ($\bar{g}_{\mu\nu}$), surrounding spatial vibrational waves ($h_{\mu\nu}^{vib}$) are drawn into the gravity well, becoming severely compressed and distorted. These compressed waves undergo extreme constructive interference, explosively increasing local tensor energy density, which in turn acts as an effective source to further curve the macroscopic spacetime.

\section{Prediction of Micro-Perturbations in Gravitational Lensing}

\subsection{Stochastic Geodesic Deviation}
To preserve the Einstein Equivalence Principle, we do not introduce arbitrary forces into the right-hand side. Instead, because the background metric itself contains high-frequency spatial vibrations, the Christoffel symbols undergo stochastic fluctuations. While the macroscopically averaged $\tilde{V}_{\mu\nu}$ generates spacetime curvature (gravity), the microscopic high-frequency perturbations $h_{\mu\nu}^{vib}$ themselves act as the source of stochastic scattering for individual photon trajectories. Light perfectly follows the null geodesic of this perturbed metric, resulting in a \textbf{`Stochastic Geodesic Deviation'}:
\begin{equation}
\frac{d^2 x^\mu}{d\lambda^2} + \left( \bar{\Gamma}^\mu_{\rho\sigma} + \delta \Gamma^\mu_{\rho\sigma}(h^{vib}) \right) \frac{dx^\rho}{d\lambda} \frac{dx^\sigma}{d\lambda} = 0
\end{equation}

\subsection{Falsifiable Prediction: Achromatic Optical Blurring}
High-resolution observations of gravitational lenses should inevitably reveal sub-pixel scattered \textbf{`Optical Blurring'} at the edges due to stochastic geodesic deviation. Conventional astronomy might dismiss this as simple scattering by intergalactic gas or plasma. 

To establish a crucial falsification test, we emphasize the \textbf{Achromaticity} of our prediction. Plasma scattering is wavelength-dependent (chromatic aberration). Conversely, the micro-scattering predicted by our model arises from pure geometrical spacetime vibrations. Therefore, we predict a strictly \textbf{`Achromatic Blurring'}---an identical blurring width across all frequency bands from radio waves to gamma rays. This serves as a definitive falsification standard.
Additionally, a \textbf{`Surplus Phase Delay ($\Delta t_{vibration}$)'} proportional to the path integral of the spatial vibration tensor is predicted, supplementary to the classical Shapiro time delay \cite{Shapiro1964}.

\subsection{Geometrical Path-Shortening and Chronological Illusions}
If starlight from deep space penetrates a macroscopic \textbf{`Spatial Junction'}---where severely twisted space folds onto itself, short-circuiting the internal geometrical distance---light bypasses immense physical distances. 

To avoid violating causality via traversable wormholes (which require unphysical negative energy), we propose that this junction is dynamically sustained by the extreme \textbf{Negative Pressure ($p_{vib} < 0$)} generated via the chameleon mechanism (detailed in Sec. IV.A). Furthermore, this \textbf{`Geometrical Path-Shortening'} does not violate the local speed of light $c$. The tensor condensation acts as an extreme geometric magnifier. Consequently, the `impossibly mature early galaxies' recently discovered by JWST are a profound \textbf{Chronological Illusion} caused by light traveling through these geometrically shortened paths, arriving earlier than expected while strictly preserving relativistic causality.

\section{The Unified Dark Fluid: Chameleon Dynamics of the Tensor Sector}
This chapter explores an alternative to standard $\Lambda$CDM cosmology \cite{Riess1998} by interpreting Dark Matter and Dark Energy as density-dependent phase transitions of $\tilde{V}_{\mu\nu}$.

\subsection{Chameleon Mechanism: Variable Equation of State}
To avoid an arbitrary transition of the tensor's properties, we adopt a dynamic analogous to Chameleon Cosmology \cite{Khoury2004}. The equation of state parameter $w = p_{vib} / (\rho_{vib} c^2)$ of the vibration tensor dynamically couples to the ambient baryonic matter density ($\rho_m$). We propose a heuristic transition function governing this phase shift:
\begin{equation}
w(\rho_m) \approx -1 + \frac{1}{1 + e^{-(\rho_m - \rho_{crit})/\sigma}}
\end{equation}
where $\sigma$ is a scaling parameter (with the same physical dimensions as density $\rho_m$) governing the sharpness of the phase transition. When the matter density drops below a critical threshold ($\rho_{crit}$), the kinetic energy is suppressed, and potential tension dominates, smoothly converging the tensor fluid to a dark energy state ($w \to -1$).

\subsection{Dark Matter as Tensor Condensation ($w \approx 0$)}
In galactic halos ($\rho_m > \rho_{crit}$), `Scale Inversion' allows spatial vibration energy to regain dynamic relevance. When fierce vibrational waves are trapped within the galactic well, they form highly coherent, rapidly oscillating `standing waves' (analogous to a macroscopic Bose-Einstein Condensate). In this phase-locked state, the rapid oscillations time-average the effective kinematic pressure to zero ($\langle p_{vib} \rangle \approx 0$), perfectly mimicking the equation of state of cold dark matter ($w \approx 0$).

This surge in energy density is defined as \textbf{`Tensor Condensation'}. By mass-energy equivalence, this assumes an `Effective Mass' \cite{Verlinde2017}. To rigorously account for the curvature of spacetime assuming spherical symmetry, the effective surplus mass is calculated using the proper radial metric element ($\sqrt{g_{rr}}$):
\begin{equation}
M_{eff}(r) = M_{visible}(r) + \frac{4\pi}{c^2} \int_{0}^{r} \langle \tilde{V}_{00} \rangle \sqrt{g_{rr}} \, r'^2 dr'
\end{equation}
Dark matter is thus phenomenologically reinterpreted as the geometrical surplus gravity generated by condensed spatial vibration energy boiling within the galactic gravity well.

\subsection{Dark Energy as Vibrational Radiation Pressure ($w \to -1$)}
In vast Cosmic Voids ($\rho_m \ll \rho_{crit}$), the chameleon coupling shifts to $w \to -1$. The intrinsic geometrical vibration of space fiercely radiates a powerful repulsive pressure outward:
\begin{equation}
p_{vib} \approx -\rho_{vib} c^2
\end{equation}
Substituting this into the Friedmann acceleration equation yields:
\begin{equation}
\frac{\ddot{a}}{a} = -\frac{4\pi G}{3} \left( \rho_m + \mathbf{\rho_{vib}} + \frac{3p_{vib}}{c^2} \right)
\end{equation}
As the universe expands and matter dilutes, in vast Cosmic Voids ($\rho_m \ll \rho_{crit}$), where $p_{vib} \approx -\rho_{vib} c^2$, the dark fluid term becomes $(\rho_{vib} - 3\rho_{vib}) = -2\rho_{vib}$. This negative effective active gravitational mass overwhelmingly drives a deterministic positive feedback loop of accelerated expansion ($\ddot{a}>0$).

\section{Conclusion}
By coupling the Spatial Vibration Tensor with the Einstein field equations via Isaacson averaging and introducing a chameleon-like dynamic equation of state, this paper provides a unified phenomenological reinterpretation of Dark Matter and Dark Energy without violating the equivalence principle or local energy conservation. The predictions of achromatic optical blurring and chronological illusions caused by geometrical path-shortening elevate this model to a testable cosmological framework.

\paragraph{Future Outlook: Cosmological Spatial Plate Tectonics.}
Assuming the expanding universe cannot maintain perfect global phase coherence, it must have inevitably fragmented into giant \textbf{`Topological Domains'} possessing distinct intrinsic vibrational phases. The final sequel, \textbf{[Mechanics of Spatial Vibration III]}, will trace the wave friction at the macroscopic boundaries (Fault Lines) of these domains. To address the Domain Wall Problem, we will discuss the early dilution of these topological defects during inflation, and investigate how their ruptures trigger Fast Radio Bursts (FRBs) via gravitational-electromagnetic resonance (Gertsenshtein effect).

\section*{Acknowledgments}
\textbf{Note on Authorship:} This work was authored by Kwang Yong Yoo. The author acknowledges the use of AI-assisted tools (Gemini) for language refinement and simulation development support. This study is an independent research work conducted under the author's sole academic responsibility.

\begin{thebibliography}{99}
\bibitem{Einstein1916} A. Einstein, Die Grundlage der allgemeinen Relativit\"{a}tstheorie, \textit{Ann. Phys.} \textbf{354}, 769 (1916).
\bibitem{Isaacson1968} R. A. Isaacson, Gravitational Radiation in the Limit of High Frequency. II. Nonlinear Terms and the Effective Stress Tensor, \textit{Phys. Rev.} \textbf{166}, 1272 (1968).
\bibitem{Shapiro1964} I. I. Shapiro, Fourth Test of General Relativity, \textit{Phys. Rev. Lett.} \textbf{13}, 789 (1964).
\bibitem{Rubin1980} V. C. Rubin \textit{et al.}, Rotational Properties of 21 SC Galaxies with a Large Range of Luminosities and Radii, \textit{Astrophys. J.} \textbf{238}, 471 (1980).
\bibitem{Khoury2004} J. Khoury, and A. Weltman, Chameleon Cosmology, \textit{Phys. Rev. D} \textbf{69}, 044026 (2004).
\bibitem{Verlinde2017} E. P. Verlinde, Emergent Gravity and the Dark Universe, \textit{SciPost Phys.} \textbf{2}, 016 (2017).
\bibitem{Riess1998} A. G. Riess \textit{et al.}, Observational Evidence from Supernovae for an Accelerating Universe and a Cosmological Constant, \textit{Astron. J.} \textbf{116}, 1009 (1998).
\end{thebibliography}

\end{document}